\expandafter\def\csname HistoricalNoteText8\endcsname{Als ich vor einiger Zeit die persönliche Bekanntschaft des Hrn Prof.\ \textsc{Gauss} zu machen das Glück hatte, sah ich unter dessen Papieren den hier folgenden schon vor mehreren Jahren entworfenen und noch nirgends bekannt gemachten Aufsatz, der die frühere Methode des Verfassers zu Bestimmung der Planetenbahnen enthält. Da ich mich bei der flüchtigen Durchsicht dieser summarischen Übersicht bald überzeugte, dass die hier von dem Verfasser entwickelte Methode zu erster genäherter Bestimmung zweier Distanzen des Planeten von der Erde wesentlich von der verschieden sei, die der Verfasser nun in seinem grössern Werk öffentlich dargelegt hat, so bat ich ihn um die Erlaubniss, diesen Aufsatz bekannt machen zu dürfen, in der Voraussetzung, dass es allen Kennern interessant sein muss, die verschiedenen Wege zu kennen, auf denen es dem Verfasser gelungen ist, zu der vollendeten Auflösung zu gelangen, von der wir unsern Lesern im vorigen Hefte eine Übersicht mitgetheilt haben. Ich hatte anfangs die Absicht, den Aufsatz mit einigen Bemerkungen zum Behuf einer Vergleichung der frühern und spätern Methode des Verfassers zu begleiten; allein da diese, hätten sie wirklich erläuternd sein sollen, etwas weitläufig und ohne Hinweisung auf das Werk selbst doch immer undeutlich geblieben wären, so schien es mir zweckmässiger, den ganzen Aufsatz (der denn doch mehr für Kenner bestimmt ist, die das Werk selbst dabei zur Hand haben), so wie er vor sechs Jahren vom Verfasser niedergeschrieben wurde, ohne allen fernern Beisatz, den astronomischen Lesern dieser Zeitschrift mitzutheilen.\hfill\textsc{von Lindenau.}}
\clearpage\SourceLeaf{153}%
\pdfbookmark[1]{Summarische Übersicht der zur Bestimmung der Bahnen der beiden neuen Hauptplaneten angewandten Methoden}{summarische1809}
\fancyhead[L]{\small SUMMARISCHE ÜBERSICHT}
\begin{center}
{\small SUMMARISCHE ÜBERSICHT DER ZUR BESTIMMUNG DER BAHNEN DER BEIDEN NEUEN HAUPTPLANETEN ANGEWANDTEN METHODEN\HistoricalNote{8}.\par}
\medskip\hrule\medskip
{\small Monatliche Correspondenz zur Beförderung der Erd- und Himmels-Kunde,\par
herausgegeben vom Freiherrn \textsc{von Zach}. September 1809.\par}
\medskip\hrule
\end{center}\medskip
\begin{center}1.\end{center}
Die von Kreis- und Parabel-Hypothesen unabhängige Bestimmung der Bahn eines Himmelskörpers aus einer kurzen Reihe von Beobachtungen beruht auf zwei Forderungen: I. Muss man Mittel haben, die Bahn zu finden, die drei gegebenen vollständigen Beobachtungen Genüge thut. II. Muss man die so gefundene Bahn so verbessern können, dass die Differenzen der Rechnung von dem ganzen Vorrath der Beobachtungen so gering als möglich werden.

\SourceLeaf{154}%
Die bequemste Art der II\textsuperscript{ten} Forderung Genüge zu leisten, scheint ihre Zurückführung auf die I\textsuperscript{ste} zu sein. Es seien für die Zeiten $t$, $t'$, $t''$ u.\ s.\ w.\ die beobachteten Orte $m$, $m'$, $m''$ (deren jeder zweifach sein wird); die nach bekannten Elementen $e$ (sechsfach anzusehn) berechneten Örter $p$, $p'$, $p''$ u.\ s.\ w., endlich die nach (noch als unbestimmt anzusehenden) Elementen $f$ berechneten Örter $q$, $q'$, $q''$ u.\ s.\ w.\ Die Differenzen der Elemente $e$ sind also
\[
p-m,\quad p'-m',\quad p''-m''\ \text{u.\ s.\ w.}
\]
Die Differenzen der Elemente $f$ hingegen
\[
(p-m)+(q-p),\quad(p'-m')+(q'-p'),\quad(p''-m'')+(q''-p'')\ \text{u.\ s.\ w.}
\]
Diese letztern sollen nun so klein als möglich werden und keine Regularität behalten.

Die Differenzen $q-p$, $q'-p'$ u.\ s.\ w.\ sind, in so fern man die Elemente $f$ als beständig ansieht, Functionen der Zeit, und da sie der Natur der Sache nach an sich klein sein werden, so darf man bei der kurzen Dauer der Beobachtungen voraussetzen, dass sie, wenn man sie für zwei äussere und eine mittlere als gegeben annimmt, für die zwischenliegenden Zeiten hinreichend genau durch Interpolation gefunden werden. Man bezeichne sie für jene drei Zeiten durch $x$, $y$, $z$ (jeder als zweifach anzusehen), so werden sie nach bekannten Gründen der Interpolationstheorie eine linearische Form haben $\alpha x+\beta y+\gamma z$, wo die Coëfficienten $\alpha$, $\beta$, $\gamma$ von der Zeit abhängen. Diese Differenzen der Elemente $f$ werden also für die Zeiten $t$, $t'$, $t''$ u.\ s.\ w.\ die Gestalt haben:
\[
\begin{gathered}
p-m+\alpha x+\beta y+\gamma z\\
p'-m'+\alpha'x+\beta'y+\gamma'z\\
p''-m''+\alpha''x+\beta''y+\gamma''z\ \text{u.\ s.\ w.}
\end{gathered}
\]
wo alles ausser $x$, $y$, $z$ bekannt ist. Man wird alsdann leicht beurtheilen können, welches die zweckmässigsten Werthe für $x$, $y$, $z$ sind. \emph{Es lässt sich zwar eine ganz methodische Anweisung geben, diese Werthe durch Rechnung zu finden;} allein ein gewisser Tact wird immer eben so sicher leiten.

Da also offenbar, so bald $x$, $y$, $z$ bestimmt sind, die II\textsuperscript{te} Forderung auf die erste zurückgeführt ist, so können wir uns blos auf diese einschränken.

\SourceLeaf{155}%
\begin{center}\emph{Bestimmung der Bahn aus drei vollständigen Beobachtungen.}\end{center}
\begin{center}2.\end{center}
Es würde zwar nicht schwer fallen, die Relation der sechs unbekannten Grössen zu den gegebenen in sechs Gleichungen darzustellen. Allein da dieselben viel zu unbehülflich ausfallen, um im geringsten brauchbar zu sein, so muss man sich begnügen, stufenweise zu derjenigen Bahn zu gelangen, die die drei Beobachtungen genau darstellt. Offenbar müssen alle überhaupt hiezu dienliche Methoden am Ende einerlei Resultat geben; die Güte des Endresultats ist also kein Maassstab für den Werth der Methode, sondern nur für die Schärfe der zum Grunde gelegten Beobachtungen. Der Werth der Methode kann nur nach der Anzahl und Bequemlichkeit der Stufen geschätzt werden, und eine Methode, wodurch man zu einer genauen Darstellung der drei Beobachtungen nicht allezeit nach Gefallen gelangen könnte, würde nicht eine schlechtere, sondern gar keine Methode sein.

Die Untersuchung zerfällt also in zwei Theile, eine \emph{erste Annäherung} und die \emph{Correctionsmethoden}. Jene wird sich auf gewisse aus der Natur des Problems geschöpfte, beinahe wahre Relationen gründen, welche von der Art sind, dass sie desto weniger fehlen, je näher die Beobachtungen einander liegen, und, mathematisch zu reden, für unendlich nahe Beobachtungen streng genau sind. Man weicht daher allerdings dem Einflusse ihrer Abweichung von der Wahrheit desto mehr aus, je nähere Beobachtungen man zum Grunde legt, wodurch die Correctionsmethoden desto mehr erleichtert oder gar entbehrlich gemacht werden. Allein dabei hat man zu erwägen, dass bei nahen Beobachtungen geringe Fehler der Beobachtungen sehr stark, zuweilen enorm auf die Elemente wirken, und daher die nachmalige Verbesserung nach der ganzen Beobachtungsreihe, die oben mit dem Namen der zweiten Forderung bezeichnet ist, desto schwieriger ausfällt. Allgemeine Regeln lassen sich über die zweckmässigste Auswahl der Beobachtungen nicht wohl geben. Es ist dabei die Güte derselben, ihre mehr oder weniger vortheilhafte Lage, die Art der Bahn selbst in Betracht zu ziehen. Bei der Ceres liessen sich mit gutem Erfolg sogleich die äussersten 41 Tage entfernten Beobachtungen zur ersten Annäherung anwenden, und die Correctionsrechnungen waren gar nicht beschwerlich. Auch bei Berechnung der II\textsuperscript{ten} Pallasbahn wurden die vorhergegangenen Näherungen nicht benutzt, sondern die erste
\SourceLeaf{156}%
Annäherungsmethode sogleich von neuem auf die 27tägigen Beobachtungen angewandt. Bei Bahnen, die der Parabel näher kämen, und wo die geocentrische Bewegung sehr schnell ist, würde man wohl lieber mit etwas kürzern Zwischenzeiten den Anfang machen. Eine durch Erfahrung geübte Beurtheilung ist hier die beste Führerin.

\begin{center}
\emph{Hauptmomente der ersten Annäherung.}\par\medskip
Erstes Moment.\par\smallskip
\emph{Genäherte Bestimmung der Abstände von der Erde in den beiden äussern Beobachtungen.}
\end{center}
\begin{center}3.\end{center}
Um bei der grossen Anzahl der in folgender Untersuchung nöthigen Zeichen die Übersicht zu erleichtern, sollen analoge Dinge bei der Erde $P$ und dem beobachteten Planeten $p$ durchaus mit einerlei Charakteren angedeutet werden, nur bei jener mit grossen, bei diesem mit kleinen. Kommt einerlei Buchstabe sowohl ohne Accent als mit einem und zweien vor, so hat man vorauszusetzen, dass der zweite und dritte eine ähnliche Beziehung auf eine zweite und dritte Beobachtung für die Zeiten $\tau'$, $\tau''$ haben, wie der erste auf die Beobachtung von der Zeit $\tau$. Übrigens ist an sich nicht nöthig, dass die Zeit $\tau'$ zwischen $\tau$ und $\tau''$ falle; inzwischen ist der Gebrauch der folgenden Vorschriften am vortheilhaftesten, wenn $\tau'$ ungefähr zwischen $\tau$ und $\tau''$ in der Mitte liegt.

$S$ Ort der Sonne (im Raum) als unbeweglich betrachtet. $p$ Ort des Planeten $p$ zur Zeit $\tau$. Hieraus erklären sich $p'$, $p''$, $P$, $P'$, $P''$.

$\mathfrak{X}$, $\mathfrak{Y}$, $\mathfrak{Z}$ drei feste willkürliche Ebenen, die einander im Mittelpunkt der Sonne senkrecht schneiden.

$x$, $y$, $z$ die senkrechten Abstände des Planeten $p$ von diesen drei Ebenen zur Zeit $\tau$. Hieraus $x'$, $y'$, $z'$; $x''$, $y''$, $z''$; $X$, $Y$, $Z$; $X'$, $Y'$, $Z'$; $X''$, $Y''$, $Z''$.
\[
\left.\begin{aligned}
\xi&=x-X\\
\eta&=y-Y\\
\zeta&=z-Z
\end{aligned}\right|\quad\text{Hieraus }\xi',\ \eta',\ \zeta';\ \xi'',\ \eta'',\ \zeta''
\]

Es sind also $\xi$, $\eta$, $\zeta$ die senkrechten Abstände des Planeten $p$ von drei beweglichen mit $\mathfrak{X}$, $\mathfrak{Y}$, $\mathfrak{Z}$ parallel durch $P$ gelegten Ebenen.

\SourceLeaf{157}%
\[
\begin{array}{c|c|c|c|c|l}
r&&\text{des }p&&S&\text{alle positiv.}\\
R&\text{Abstände}&\text{des }P&\text{von}&S&\text{Hieraus }r'\ \text{u.\ s.\ w.}\\
\rho&&\text{des }p&&P&
\end{array}
\]
\[
\begin{array}{c|c|c|c|l}
b&&Sp&&\mathfrak{Z}\\
B&\text{Winkel der Linie}&SP&\text{mit der Ebene}&\mathfrak{Z}\\
\beta&&Pp&&\text{die }\mathfrak{Z}\text{ parallel ist}
\end{array}
\]
\[
\left.\begin{aligned}
d&=r\cos b\\D&=R\cos B\\\delta&=\rho\cos\beta
\end{aligned}\right|\quad
\begin{gathered}\text{d.\ i.\ projicirte Abstände auf die Ebene }\mathfrak{Z}\\\text{und die ihr parallele}\end{gathered}
\]
\[
\begin{array}{c|c|l}
l&&\mathfrak{Y}\\
L&\text{Winkel dieser Projection mit der Ebene}&\mathfrak{Y}\\
\lambda&&\text{die }\mathfrak{Y}\text{ parallel.}
\end{array}
\]
Die Winkel $b$ und $l$ sind auf eben den Seiten von $\mathfrak{Z}$ und $\mathfrak{Y}$ als positiv anzunehmen, auf welchen man $z$ und $y$ als positiv ansieht. Die Winkel $b$ kann man immer zwischen den Grenzen $-90^0$, $+90^0$ lassen (dass die $d$ u.\ s.\ w.\ immer positiv bleiben); die Winkel $l$ hingegen kann man von $0$ bis $360^0$ wachsen lassen, und zwar so, dass man sie da $\left\{\begin{smallmatrix}0\\180^0\end{smallmatrix}\right\}$ setzt, wo die $x$ $\left\{\begin{smallmatrix}\text{positiv}\\\text{negativ}\end{smallmatrix}\right\}$ sind. Auf die Weise hat man
\[
\left.\begin{aligned}
x&=r\cos b\cos l=d\cos l\\
y&=r\cos b\sin l=d\sin l\\
z&=r\sin b\phantom{\cos l}=d\operatorname{tang}b
\end{aligned}\right|\quad
\begin{gathered}\text{und ähnliche Gleichungen für}\\x'\ \text{u.\ s.\ w.}\quad X\ \text{u.\ s.\ w.}\quad\xi\ \text{u.\ s.\ w.}\end{gathered}
\]

Die Bahnen von $p$ und $P$ nehmen wir in Ebenen an, indem wir von den fremden Einwirkungen, wodurch sie afficirt werden, abstrahiren. Wir setzen Länge von $p$ in der Bahn zur Zeit $\tau$, $=v$ (also $v'$, $v''$; $V$, $V'$, $V''$); und machen $\frac12r'r''\sin(v''-v')=f$, $\frac12r''r\sin(v-v'')=f'$, $\frac12rr'\sin(v'-v)=f''$. Es sind folglich $f$, $-f'$, $f''$ die Flächen der Dreiecke $p'Sp''$, $pSp''$, $pSp'$ positiv (vorausgesetzt, dass $p$ rechtläufig ist und $\tau'$ zwischen $\tau$ und $\tau''$ liegt; das Arrangement der Zeichen für andere Fälle hat keine Schwierigkeit). Auf ähnliche Art $F$, $F'$, $F''$. Durch $g$, $-g'$, $g''$, $G$, $-G'$, $G''$ bezeichnen wir die Flächen der Ausschnitte aus der ganzen Bahn, die diesen Dreiecken entsprechen, deren Zeichen wir denen von $f$, $-f'$, $f''$, $F$, $-F'$, $F''$ gleich voraussetzen. Es sind daher $g$, $g'$, $g''$ und $G$, $G'$, $G''$ den Zeitintervallen $\tau''-\tau'$, $\tau-\tau''$, $\tau'-\tau$ proportional.

\SourceLeaf{158}%
Die Bahnen von $p$ und $P$ sind Kegelschnitte, deren halbe grosse Achsen wir mit $a$, $A$ bezeichnen. Die Excentricität der Bahn von $p$ machen wir $=e=\sin\varphi$ (für eine Ellipse); daher $a\cos\varphi^2=k$ der halbe Parameter wird. Länge der Sonnenferne von $p$ in der Bahn $=\pi$. Mittlere Länge $=m$ (daher $m'$, $m''$, $M$, $M'$, $M''$). Andere Zeichen sollen im Verfolg der Untersuchung selbst angezeigt werden.

\begin{center}4.\end{center}
Daraus, dass $p$, $p'$, $p''$ mit $S$ in einer Ebene sind, folgt nach einem bekannten Lehrsatz
\[
0=xy'z''+x'y''z+x''yz'-xy''z'-x'yz''-x''y'z
\]
und hieraus, dass von folgenden neun Grössen
\[
\begin{array}{ccc}
y'z''-y''z'&y''z-yz''&yz'-y'z\\
z'x''-z''x'&z''x-zx''&zx'-z'x\\
x'y''-x''y'&x''y-xy''&xy'-x'y
\end{array}
\]
die drei obern den drei mittlern und den drei untern resp.\ proportional sind. Man wird sich leicht überzeugen,

I. Dass eben diesen Grössen auch $f$, $f'$, $f''$ proportional sind, da die drei obern, mittlern, untern nur das doppelte Areal der Projection der Dreiecke, deren Inhalt $f$, $f'$, $f''$ sind, auf die Fundamental-Ebenen $\mathfrak{X}$, $\mathfrak{Y}$, $\mathfrak{Z}$, vorstellen und sich also zu diesen verhalten wie die doppelten Cosinus der Neigung der Bahn von $p$ gegen jene Ebene zur Einheit (In einer vollständigen Abhandlung würden hier noch Erinnerungen wegen der Zeichen nöthig sein, die man aber auch durch blossen Calcul leicht umgehen kann).

II. Dass, wenn man die drei obern, die drei mittlern oder die untern mit $x$, $x'$, $x''$. mit $y$, $y'$, $y''$ oder mit $z$, $z'$, $z''$ multiplicirt, die Summe der Producte $=0$ wird. Hieraus lässt sich leicht schliessen
\[
\begin{array}{r|c|l}
0=fx+f'x'+f''x''&\text{Durch ganz}&0=FX+F'X'+F''X''\\
0=fy+f'y'+f''y''&\text{analoge Schlüsse}&0=FY+F'Y'+F''Y''\\
0=fz+f'z'+f''z''&\text{hat man}&0=FZ+F'Z'+F''Z''
\end{array}
\]
\SourceLeaf{159}%
Hieraus lassen sich leicht folgende drei Gleichungen ableiten
\[
\begin{aligned}
&(F+F'')(f\xi+f'\xi'+f''\xi'')\\
&\qquad=(Ff''-F''f)(X-X'')+[F'(f+f'')-(F+F'')f']X'\\
&(F+F'')(f\eta+f'\eta'+f''\eta'')\\
&\qquad=(Ff''-F''f)(Y-Y'')+[F'(f+f'')-(F+F'')f']Y'\\
&(F+F'')(f\zeta+f'\zeta'+f''\zeta'')\\
&\qquad=(Ff''-F''f)(Z-Z'')+[F'(f+f'')-(F+F'')f']Z'
\end{aligned}
\]

Aus diesen drei Gleichungen leiten wir vier andere ab, indem wir sie respective
\[
\begin{array}{c|c|c|c|l}
\text{zuerst mit}&\text{dann mit}&\text{dann mit}&\text{endlich mit}&\text{multipli-}\\
\eta\zeta''-\eta''\zeta&\eta Z'-Y'\zeta&\eta'Z'-Y'\zeta'&\eta''Z'-Y'\zeta''&\text{ciren und}\\
\zeta\xi''-\zeta''\xi&\zeta X'-Z'\xi&\zeta'X'-Z'\xi'&\zeta''X'-Z'\xi''&\text{d.\ Producte}\\
\xi\eta''-\xi''\eta&\xi Y'-X'\eta&\xi'Y'-X'\eta'&\xi''Y'-X'\eta''&\text{addiren.}
\end{array}
\]

Zur bequemen Übersicht bezeichnen wir die Summen der Producte, die entstehen
\begin{center}\begingroup\fontsize{9}{11}\selectfont\setlength{\tabcolsep}{.8pt}
\begin{tabular}{c|c|c|c|c|c}
&&&&&\shortstack{nach d.\ Multi-\\plication dieser\\Factoren mit}\\\hline
&$*$&$*$&$\delta'\delta D'[\pi'\pi P']$&$\delta''\delta D'[\pi''\pi P']$&$\xi,\quad\eta,\quad\zeta$\\
$[\pi\pi'\pi'']$&$\delta\delta'\delta''\times[\pi\pi'\pi'']$&$\delta\delta'D'[\pi\pi'P']$&$*$&$\delta''\delta'D'[\pi''\pi'P']$&$\xi',\quad\eta',\quad\zeta'$\\
&$*$&$\delta\delta''D'[\pi\pi''P']$&$\delta'\delta''D'[\pi'\pi''P']$&$*$&$\xi'',\quad\eta'',\quad\zeta''$\\
$[\pi P\pi'']$&$\delta D\delta''[\pi P\pi'']$&$\delta DD'[\pi PP']$&$\delta'DD'[\pi'PP']$&$\delta''DD'[\pi''PP']$&$X,\quad Y,\quad Z$\\
$[\pi P'\pi'']$&$\delta D'\delta''[\pi P'\pi'']$&$*$&$*$&$*$&$X',\quad Y',\quad Z'$\\
&$\delta D''\delta''[\pi P''\pi'']$&$\delta D''D'[\pi P''P']$&$\delta'D''D'[\pi'P''P']$&$\delta''D''D'[\pi''P''P']$&$X'',\quad Y'',\quad Z''$
\end{tabular}\endgroup\end{center}
Es ist klar, dass in die hier mit $*$ ausgefüllten Stellen $0$ kommen müsse, und dass alle durch eingeklammerte Zeichen angedeuteten Grössen \emph{gegeben} sind. Nemlich
\[
\begin{aligned}\relax
[\pi\pi'\pi'']&=\operatorname{tang}\beta\sin(\lambda'-\lambda'')+\operatorname{tang}\beta'\sin(\lambda''-\lambda)+\operatorname{tang}\beta''\sin(\lambda-\lambda')\\
[\pi P\pi'']&=\operatorname{tang}\beta\sin(L-\lambda'')+\operatorname{tang}B\sin(\lambda''-\lambda)+\operatorname{tang}\beta''\sin(\lambda-L)\\
&\text{u.\ s.\ w.}
\end{aligned}
\]

Es ist nicht nöthig alle 16 Gleichungen herzusetzen, da sie alle auf analoge Art aus der ersten abgeleitet werden können, indem man nur $\beta$ mit
\expandafter\def\csname HistoricalNoteText9\endcsname{Ich erlaube mir diesen leicht verständlichen Ausdruck der Kürze wegen. Der positive Pol von $\mathfrak{X}$ liegt auf der Seite dieser Ebene, wo die $x$ positiv gezählt werden u.\ s.\ f.}%
\SourceLeaf{160}%
$\beta'$, $\beta''$, $B$, $B'$, $B''$ und $\lambda$ mit $\lambda'$, $\lambda''$, $L$, $L'$, $L''$ vertauscht, wenn an der Stelle von $\pi$ resp.\ $\pi'$, $\pi''$, $P$, $P'$, $P''$ steht u.\ s.\ f.\ Zugleich sieht man, dass die 16 Grössen sich auf 12 reduciren, da
\[
\begin{aligned}
+[\pi P'\pi'']&=-[\pi\pi''P']=+[\pi''\pi P']\\
[\pi\pi'P']&=-[\pi'\pi P']\\
[\pi''\pi'P']&=-[\pi'\pi''P']
\end{aligned}
\]

Ferner erkennt man leicht, dass der Ausdruck $[\pi\pi'\pi'']$, multiplicirt in das Product aus den drei Cosinussen der darin vorkommenden Breiten, der sechsfache Inhalt einer Pyramide ist, deren Spitze in den Mittelpunkt und die drei Winkelpunkte der Basis in die Oberfläche einer mit dem Halbmesser $1$ beschriebenen Kugel so fallen, dass sie den drei geocentrischen Örtern von $p$ entsprechen, und zwar wird jenes Zeichen den sechsfachen Werth positiv oder negativ angeben, je nachdem jene drei geocentrischen Örter auf der Sphäre in entgegengesetzter oder gleicher Ordnung (sens) liegen, wie die positiven\HistoricalNote{9} Pole der Ebenen $\mathfrak{X}$, $\mathfrak{Y}$, $\mathfrak{Z}$ resp.\ Und ganz ähnliche Dinge drücken die übrigen Zeichen aus.

Auf diese Weise entspringen folgende vier Gleichungen:
\[
\begin{aligned}
1)\quad &(F+F'')f'\delta'[\pi\pi'\pi'']\\
&=(Ff''-F''f)(D[\pi P\pi'']-D''[\pi P''\pi''])\\
&\qquad+(F'(f+f'')-(F+F'')f')D'[\pi P'\pi'']\\
2)\quad &(F+F'')(f'\delta'[\pi\pi'P']+f''\delta''[\pi\pi''P'])\\
&=(Ff''-F''f)(D[\pi PP']-D''[\pi P''P'])\\
3)\quad &(F+F'')(f\delta[\pi'\pi P']+f''\delta''[\pi'\pi''P'])\\
&=(Ff''-F''f)(D[\pi'PP']-D''[\pi'P''P'])\\
4)\quad &(F+F'')(f\delta[\pi''\pi P']+f'\delta'[\pi''\pi'P'])\\
&=(Ff''-F''f)(D[\pi''PP']-D''[\pi''P''P'])
\end{aligned}
\]

\begin{center}5.\end{center}
Wir wollen nun diese vier Gleichungen, die streng richtig sind, näher betrachten, um darauf die erste Annäherung zu gründen. Betrachten wir die Zwischenzeiten als unendlich kleine Grössen der ersten Ordnung, so sind $f$, $f'$, $f''$, $G$, $G'$, $G''$ und alle eingeklammerte Grössen von derselben ersten Ordnung mit Ausnahme von $[\pi\pi'\pi'']$, welches von der dritten Ordnung ist. Die Beweise, so wie die sich leicht darbietenden Bemerkungen über specielle Ausnah%
\SourceLeaf{161}%
men, übergehe ich. Müsste man die Neigung der Bahnen von $p$ und $P$ gegen einander als Grössen der ersten Ordnung betrachten, so würden alle eingeklammerten Grössen eine Ordnung höher stehen. Ferner ist
\[
Ff''-F''f=\frac FG.\frac{f''}{g''}Gg''-\frac{F''}{G''}.\frac fg G''g
\]
oder (weil $Gg''=G''g$)
\[
=\left(\frac FG.\frac{f''}{g''}-\frac{F''}{G''}.\frac fg\right)Gg''
\]
Nun ist $G-F$ eine Grösse der dritten Ordnung, daher $1-\frac FG$ eine der zweiten u.\ s.\ w.\ und folglich auch $\frac FG.\frac{f''}{g''}-\frac{F''}{G''}.\frac fg$ eine der zweiten und mithin $Ff''-F''f$ von der vierten (es würde sogar von der fünften sein, wenn $\tau'$ mitten zwischen $\tau$ und $\tau''$ fiele). Was daher oben in der zweiten, dritten und vierten Gleichung rechts steht, ist von der fünften Ordnung, von dem, was links steht, ist sowohl der erste als der zweite Theil von der dritten, man kann also zur ersten Annäherung setzen
\[
\begin{array}{ll}
\text{aus 2)}&f'\delta'[\pi\pi'P']=-f''\delta''[\pi\pi''P']\\
\text{aus 4)}&f\delta[\pi\pi''P']=-f'\delta'[\pi'\pi''P']
\end{array}
\]
Was aus 3) folgt, ist mit diesen Resultaten identisch. Zur fernern Abkürzung können wir hier noch $\frac fg=\frac{f''}{g''}$ setzen, welche Grössen beide von der Einheit nur um die zweite Ordnung, und eben so viel von einander verschieden sind (fällt $\tau'$ mitten zwischen $\tau$ und $\tau''$, so ist die Differenz nur von der dritten). Da nun: $g:-g':g''=\tau''-\tau':\tau''-\tau:\tau'-\tau$, so wird
\[
\begin{aligned}
5)\qquad\delta&=\frac gf.\frac{f'}{g'}.\frac{\tau''-\tau}{\tau''-\tau'}.\frac{[\pi'\pi''P']}{[\pi\pi''P']}.\delta'\\
6)\qquad\delta''&=\frac{g''}{f''}.\frac{f'}{g'}.\frac{\tau''-\tau}{\tau'-\tau}.\frac{[\pi\pi'P']}{[\pi\pi''P']}.\delta'
\end{aligned}
\]

Diese Formeln geben $\delta$ und $\delta''$ aus $\delta'$ bis auf die zweite Ordnung richtig inclusive, wenn $\tau'$ mitten zwischen $\tau$ und $\tau''$ liegt, sonst exclusive. Im letzten Falle kann man $\frac{f'}{g'}=1$ setzen, da der Unterschied nur von der zweiten Ordnung ist; im ersten hingegen ist es der Mühe nicht unwerth
\[
f'=g'+\tfrac43(f+f'+f'')\quad\text{oder}\quad\frac{f'}{g'}=1+\tfrac43\frac{f+f'+f''}{f'}
\]
zu setzen, welches bald wird näher bestimmt werden und um 1 Ordnung genauer ist (Man sieht leicht, dass $f+f'+f''$ dem Dreiecke zwischen den drei Örtern
\SourceLeaf{162}%
\expandafter\def\csname HistoricalNoteText10\endcsname{Abhandlung über die leichteste Methode u.\ s.\ w.\ S.\ 45.}%
$p$, $p'$, $p''$ gleich ist, also nach einer bekannten Annäherung $=\frac34\times$ Abschnitt der krummen Fläche zwischen der Sehne $pp''$ und dem Bogen). Übrigens folgt aus obigen Formeln
\[
\frac{\delta''}{\delta}=\frac{[\pi\pi'P']}{[\pi'\pi''P']}.\frac{\tau''-\tau'}{\tau'-\tau}
\]
welches, wenn man für $\mathfrak{Z}$ die Ekliptik annimmt oder $B$, $B'$, $B''=0$ setzt, sich in
\[
\frac{\delta''}{\delta}=\frac{\operatorname{tang}\beta\sin(\lambda'-L')-\operatorname{tang}\beta'\sin(\lambda-L')}{\operatorname{tang}\beta'\sin(\lambda''-L')-\operatorname{tang}\beta''\sin(\lambda'-L')}.\frac{\tau''-\tau'}{\tau'-\tau}
\]
d.\ i.\ in die bekannte \textsc{Olbers}sche Formel\HistoricalNote{10} verwandelt.

\begin{center}6.\end{center}
Nachdem wir aus den Formeln 2.\ 3.\ 4.\ geschmeidige Näherungen abgeleitet haben, nehmen wir auf ähnliche Weise die erste vor. Bekanntlich ist
\[
\begin{aligned}
\frac1r&=\frac1k(1-e\cos(v-\pi))\\
\frac1{r'}&=\frac1k(1-e\cos(v'-\pi))\\
\frac1{r''}&=\frac1k(1-e\cos(v''-\pi))
\end{aligned}
\]

Hieraus folgt, wenn man mit $\sin(v''-v')$, $\sin(v-v'')$, $\sin(v'-v)$ multiplicirt und addirt:
\[
\begin{aligned}
\frac{f+f'+f''}{rr'r''}&=\frac1k(\sin(v''-v')+\sin(v-v'')+\sin(v'-v))\\
&=-\frac4k\sin\tfrac12(v''-v')\sin\tfrac12(v-v'')\sin\tfrac12(v'-v)
\end{aligned}
\]
oder
\[
\frac{f+f'+f''}{f'}=-\frac{2r'}k.\frac{\sin\tfrac12(v''-v')\sin\tfrac12(v'-v)}{\cos\tfrac12(v''-v)}
\]
Nach einem bekannten Lehrsatze aus der theorischen Astronomie ist
\[
\frac{\text{Beschriebener Raum}}{\text{Mittlere Bewegung}}=\frac{a^{\frac32}.\BareRoot{k}}2
\]
Mithin
\[
k=\frac{4gg''}{a^3(m'-m)(m''-m')}=\frac{4gg''}{A^3(M'-M)(M''-M')}
\]
\SourceLeaf{163}%
\expandafter\def\csname HistoricalNoteText11\endcsname{Hier ist nemlich wirkliche Subtraction, da $[\pi P\pi'']$, $[\pi P''\pi'']$ einerlei Zeichen haben; dies ist bei den Coëfficienten von $Ff''-F''f$ in den Gleichungen 2.\ 3.\ 4.\ nicht der Fall, sondern die Theile werden da eigentlich addirt. Eine tiefere Untersuchung wäre hier zu weitläuftig.}%
Also
\[
\begin{aligned}
\frac{f+f'+f''}{f'}&=-\frac{A^3(M'-M)(M''-M)}{2\cos^2\tfrac12(v''-v)}.\frac1{r'^3}.\frac{r'r'}{rr''}\\&\quad.\frac{r'r''\sin\tfrac12(v''-v')}{g}.\frac{rr'\sin\tfrac12(v'-v)}{g''}\\
\frac{f+f'+f''}{g'}&=-\frac{r'r''\sin\tfrac12(v''-v')}{g}.\frac{rr''\sin\tfrac12(v''-v)}{g'}.\frac{rr'\sin\tfrac12(v'-v)}{g''}\\&\quad.\frac{A^3(M'-M)(M''-M)}{2r'^3}.\frac{r'r'}{rr''}
\end{aligned}
\]
Man folgert hieraus leicht, da
\[
\frac1{\cos\tfrac12(v''-v)},\qquad\frac{r'r''\sin\tfrac12(v''-v')}{g},\qquad\frac{rr'\sin\tfrac12(v'-v)}{g'}
\]
und, wenn entweder $\tau'$ in die Mitte von $\tau$ und $\tau''$ fällt oder man den Unterschied der Bahn des $p$ vom Kreise als von der ersten Ordnung betrachten kann, auch $\frac{r'r'}{rr''}$ von der Einheit nur um Grössen der zweiten Ordnung abweichen, dass man näherungsweise setzen dürfe
\[
\frac{f+f'+f''}{f'}=-\frac{A^3}{2r'^3}(M'-M)(M''-M')
\]
Auf gleiche Weise ist näherungsweise
\[
\frac{F+F'+F''}{F'}=-\frac{A^3}{2R'^3}(M'-M)(M''-M')
\]

Die letztere Grösse kann man, wenn man will, auch genau berechnen, da alles dazu gegeben ist. Beide sind von der zweiten Ordnung und bis auf die vierte excl.\ bestimmt. — Wir haben also
\[
\begin{aligned}
F'(f+f'')-(F+F'')f'&=F'(f+f'+f'')-(F+F'+F'')f'\\
&=F'f'\tfrac12A^3(M'-M)(M''-M')\left(\frac1{R'^3}-\frac1{r'^3}\right)
\end{aligned}
\]
Grösse der vierten Ordnung bis auf die sechste excl.\ bestimmt. In der Gleichung 1 oben, ist der Theil linker Hand von der fünften Ordnung; von dem Theile rechter Hand ist das erste Glied von der sechsten oder siebten Ordnung, nemlich $Ff''-F''f$ ist von der vierten oder fünften und $D[\pi P\pi'']-D''[\pi P''\pi'']$ ist von der zweiten\HistoricalNote{11}, das zweite von der fünften, wir lassen also jenes weg und bekommen dadurch
\[
7)\qquad\frac{[\pi\pi'\pi'']}{[\pi P'\pi'']}.\frac2{A^3(M'-M)(M''-M')}=\left(\frac1{R'^3}-\frac1{r'^3}\right)\frac R{\delta'}
\]
\SourceLeaf{164}%
bis auf Grössen der zweiten excl.\ richtig, wenn $\tau'$ zwischen $\tau$ und $\tau''$ in die Mitte fällt; sonst nur um Grössen der ersten unrichtig. Diese Formel, welche folgende Gestalt erhält, wenn wir für $\mathfrak{Z}$ die Ekliptik nehmen, ist der wichtigste Theil der ganzen Methode und ihre erste Grundlage.
\[
\begin{aligned}
\left\{1-\left(\frac{R'}{r'}\right)^3\right\}.\frac{R'}{\delta'}&=\frac2{A^3.(M'-M)(M''-M')}\\
&\quad.\frac{\operatorname{tg}\beta'\sin(\lambda''-\lambda)-\operatorname{tg}\beta\sin(\lambda''-\lambda')-\operatorname{tg}\beta''\sin(\lambda'-\lambda)}{\operatorname{tg}\beta\sin(L'-\lambda'')-\operatorname{tg}\beta''\sin(L'-\lambda)}
\end{aligned}
\]
$L'$ ist Länge der Sonne $+180^0$.

Da das, was hier rechts steht, gegeben ist, so sieht man, dass aus der Verbindung dieser Gleichung mit folgender
\[
\frac{\frac{R'}{\delta'}}{\frac{R'}{r'}}=\BareRoot{(1+\operatorname{tang}\beta'^2+\frac{R'R'}{\delta'\delta'}+2\frac{R'}{\delta'}\cos(\lambda'-L'))}
\]
sich $r'$ leicht finden lassen wird. Die indirecte Methode ist hier bei weitem die bequemste, man kommt nach wenigen Versuchen, wofür sich leicht zweckmässige Vorschriften geben lassen, sehr schnell zum Ziele. Man wird dabei auch allemal sehen können, ob es mehr als einen Werth für $r'$ gibt, also mehr als eine Bahn, wodurch die Beobachtungen dargestellt werden, welches allerdings zuweilen der Fall sein kann. —

Sonst ist noch zu bemerken, dass eigentlich hiebei die Längen nicht von den beweglichen Äquinoctial-Punkten, sondern von einem festen Punkte an gezählt werden müssen; in der Anwendung ist indess der Unterschied von keiner Bedeutung. Drückt man die Zeit in Tagen aus, so hat man
\[
\log(M'-M)(M''-M')=\log(\tau'-\tau)+\log(\tau''-\tau')+6.4711352(-10)
\]
($M$ etc.\ müssen nemlich nicht in Graden, sondern in Theilen des Halbmessers ausgedrückt werden).

Hat man $\delta'$ und $r'$, so lässt sich auch $\frac{f+f'+f''}{f'}$ und sonach auch $\delta$ und $\delta''$ bestimmen. Übrigens lassen sich aus der Betrachtung der Formel 7) noch andere interessante Folgerungen ableiten, die hier übergangen werden müssen.

\begin{center}Zweites Moment.\\\emph{Genäherte Bestimmung der Elemente.}\end{center}
\begin{center}7.\end{center}
Die mittlere Beobachtung für die Zeit $\tau'$ lassen wir nun ganz weg und gebrauchen dafür die Abstände $\delta$, $\delta''$, die im vorigen Moment näherungsweise be%
\SourceLeaf{165}%
stimmt sind. Es ist klar, dass daraus nunmehr die heliocentrischen Längen, Breiten und Abstände abgeleitet werden können; hieraus die Länge des $\AscendingNode$ und Neigung der Bahn und die Längen in der Bahn. Es bleibt also noch das Problem übrig
\begin{center}\begin{tabular}{ll}
Aus zwei Längen in der Bahn $\quad.\quad.$&$v,\ v''$\\
\quad den Abständen von der Sonne $\quad.$&$r,\ r''$\\
\quad den zugehörigen Zeiten $\quad.\quad.\quad.$&$\tau,\ \tau''$
\end{tabular}\end{center}
die übrigen Elemente zu bestimmen; nemlich $a$, $e$, $\pi$ und die Epoche. Da die Relationen dieser Grössen zu den gegebenen transcendent sind, so muss man sich hier wieder an indirecte Methoden halten. Wir wollen hier drei betrachten.

\par\medskip\emph{Erste Methode.}\par
Wenn man $\pi$ als gegeben voraussetzt.

Man setze
\[
\frac{r''+r}{r''-r}\operatorname{tang}\tfrac12(v''-v)=\operatorname{tang}\zeta
\]
so ist
\[
e=\frac{\cos\zeta}{\cos\tfrac12(v''-v)\cos[\pi-\tfrac12(v+v'')-\zeta]}
\]
Am rathsamsten ist es, sodann $k$ auf doppelte Art zu berechnen
\[
k=r[1-e\cos(v-\pi)]=r''[1-e\cos(v''-\pi)]
\]
welches auch zur Prüfung der Rechnung dient. Setzt man $e=\sin\varphi$, so ist $a=\frac{k}{\cos\varphi^2}$. Aus den wahren Anomalien kann man sodann entweder nach den gewöhnlichen oder bequemer nach indirecten Methoden die excentrischen und mittlern Anomalien und Längen berechnen; hieraus und aus der durch $a$ gegebenen mittlern Bewegung erhält man eine doppelte Bestimmung der mittlern Länge für eine beliebige Epoche. Stimmen beide überein, so hat man den richtigen Werth von $\pi$ getroffen; wo nicht, so muss man mit einem etwas geänderten Werthe von $\pi$ die Rechnung wiederholen und den wahren durch Interpolation finden. ‘Rathsam ist es hiebei, die übrigen Elemente \emph{nicht} durch Interpolation, sondern durch neue Rechnung aus den corrigirten Werthen von $\pi$ zu suchen, und nicht eher aufzuhören, bis die beiden Werthe für die Epoche vollkommen übereinstimmen.’
\SourceLeaf{166}%
\par\medskip\emph{Zweite Methode.}\par
Wenn man $e$ voraussetzt.

Hier ist die Rechnung ganz dieselbe, nur muss man, da hier $\pi$ durch die Gleichung
\[
\cos[\pi-\tfrac12(v+v'')-\zeta]=\frac{\cos\zeta}{e\cos\tfrac12(v''-v)}
\]
gesucht werden muss, den wahren Werth schon beiläufig kennen, weil dem Cosinus zwei verschiedene Werthe zugehören.

Übrigens ist I der II vorzuziehen, und überhaupt sind beide Methoden nur dann zweckmässig, wenn der durchlaufene Bogen schon sehr gross ist, und man die Elemente schon beinahe kennt. Bei den ersten Annäherungen aus einer kurzen Reihe von Beobachtungen hält man sich allemal an folgende

\Needspace{6\baselineskip}\par\medskip\emph{Dritte Methode.}\par
Wenn $k$ vorausgesetzt wird.

Man hat hier
\[
\begin{aligned}
\frac{\frac1{r''}-\frac1r}{2\sin\tfrac12(v''-v)}&=\frac ek\sin(\tfrac12(v+v'')-\pi)\\
\frac{\frac2k-\frac1{r''}-\frac1r}{2\cos\tfrac12(v''-v)}&=\frac ek\cos(\tfrac12(v+v'')-\pi)
\end{aligned}
\]
Die Division gibt also $\operatorname{tang}[\tfrac12(v+v'')-\pi]$, hieraus $\pi$, und darnach aus einer von beiden Gleichungen $e$. Das übrige ist ganz wie bei den vorhergehenden Methoden.

Der Vorzug dieser dritten Methode besteht darin, dass man für $k$ sogleich einen sehr genäherten Werth finden kann, wenn der Bogen $v''-v$ nicht zu gross ist. Es ist nemlich der Ausschnitt zwischen den beiden Radiis Vectoribus d.\ i.
\[
g'=\frac{a^{\frac32}\BareRoot{k}}2(m''-m)=\tfrac12A^{\frac32}(M''-M)\BareRoot{k}
\]
Nun ist $2g'=\int rr\,\mathrm{d}w$ von $w=v$ bis $w=v''$.

Nun ist aber nach der bekannten \textsc{Cotes}ischen Näherungs-Integrations-Formel $\int\varphi w.\mathrm{d}w$ von $w=v$ bis $w=v''$
\[
=(\tfrac12\varphi v+\tfrac12\varphi v'')(v''-v)
\]
\SourceLeaf{167}%
und noch genauer
\[
=(\tfrac16\varphi v+\tfrac23\varphi\tfrac12(v+v'')+\tfrac16\varphi v'')(v''-v)
\]
noch genauer
\[
=(\tfrac18\varphi v+\tfrac38\varphi(\tfrac23v+\tfrac13v'')+\tfrac38\varphi(\tfrac13v+\tfrac23v'')+\tfrac18\varphi v'')(v''-v)\quad\text{u.\ s.\ w.}
\]
Es ist aber allemal hinreichend bei den ersten beiden stehen zu bleiben.

Nach der ersten hat man also
\[
2g'=\tfrac12(rr+r''r'')(v''-v)\quad\text{und}\quad\BareRoot{k}=\frac{\tfrac12(rr+r''r'')}{A^{\frac32}}.\frac{v''-v}{M''-M}
\]
$A$ macht man gewöhnlich $=1$; $v''-v$ und $M''-M$ werden in Secunden ausgedrückt, so ist $\log(M''-M)=\log(\tau''-\tau)+3.5500073$. Zur Abkürzung der Rechnung macht man $\frac{r'}r=\operatorname{tang}(45^0\pm\psi)$, wodurch $\tfrac12(rr+r''r'')=\frac1{\cos2\psi}$.

Nach der zweiten Integrations-Formel setze man den Radius Vector, der der Länge $\tfrac12(v''+v)$ zugehört, $=r^*$, so ist
\[
\frac1{r^*}=\tfrac12\left(\frac1r+\frac1{r''}\right)+\left[\tfrac12\left(\frac1r+\frac1{r''}\right)-\frac1k\right]\frac{2\sin\tfrac14(v''-v)^2}{\cos\tfrac12(v''-v)}
\]
Hienach kann man $r^*$ vermittelst des ersten Werthes von $k$ bestimmen. Sodann ist
\[
2g'=\tfrac16(rr+r''r'')+\tfrac23r^*r^*
\]
also der neue Werth von $k$
\[
=k\left\{1+\frac{\tfrac23(r^*r^*-\tfrac12(rr+r''r''))}{rr+r''r''}\right\}^2
\]

In der Ausübung ist es gewöhnlich genau genug und bequemer, den Logarithmen des neuen Werthes von $k$ dadurch zu suchen, dass man den Logarithmen des ersten um $\tfrac43\log\frac{r^*r^*}{\tfrac12(rr+r''r'')}$ vermehrt. Will man mit diesem neuen Werthe von $k$ den Werth von $\frac1{r^*}$ nach obiger Formel nochmals genauer bestimmen und darnach den Werth von $k$ abermals berichtigen, so werden fast immer die zuletzt entspringenden zweifachen Bestimmungen der Epoche so gut übereinstimmen, dass gar keine neue Voraussetzung nöthig sein wird. Bei der $\text{\AstroSymbols ⚳}$ und $\text{\AstroSymbols ⚴}$ stimmten sie, da $\tau''-\tau$ doch 41 und 42 Tage war, immer auf ein paar Hunderttheile von Secunden.
\SourceLeaf{168}%
\begin{center}\emph{Verbesserungs-Methoden.}\end{center}
\begin{center}8.\end{center}
Berechnet man nach den durch die vorhergehenden Methoden gefundenen genäherten Elementen den Ort für die Zeit $\tau'$, und findet man denselben mit der Beobachtung übereinstimmend, so ist die Arbeit vollendet. Gewöhnlich wird die Übereinstimmung sehr gross sein (oft betrug der Unterschied bei meinen Rechnungen nur wenige Secunden), aber doch selten vollkommen, theils weil zum Theil nur genäherte Voraussetzungen zum Grunde liegen, theils \emph{weil selbst die Sonnen-Örter, die man dabei gebraucht, nicht elliptisch sind, sondern die kleinen Störungen mit einschliessen.} Man könnte nun zwar für die oben weggelassenen kleinen Grössen höherer Ordnungen aus den genäherten Elementen die Werthe sehr nahe bestimmen, und so die obigen Formeln und die Werthe von $\delta$, $\delta''$ darnach verbessern; allein ich bin der Meinung, dass diese Rechnung weit beschwerlicher sein würde, als eine von den folgenden Methoden.

Die allerleichteste erste Verbesserungs-Methode, auf die ich erst bei Veranlassung der $\text{\AstroSymbols ⚴}$ verfiel, und die ich dabei, da die Zwischenzeiten noch kurz waren, mit dem glücklichsten Erfolg angewandt habe, ist folgende.

Gesetzt nach den genäherten Elementen, \emph{die auf obige Weise gefunden waren}, sei der berechnete Ort für die Zeit $\tau$ in Länge $=\lambda'+\mathfrak{L}$, in Breite $=\beta'+\mathfrak{B}$, da der beobachtete $\lambda'$ und $\beta'$ ist, so dass durch Conspiration aller kleinen Unrichtigkeiten in den Voraussetzungen die Länge um $\mathfrak{L}$, die Breite um $\mathfrak{B}$ zu gross ausfällt; so berechne man ganz von neuem und \emph{ganz auf dieselbe Art} die Bahn, indem man die Beobachtungen
\[
\begin{array}{ccc}
\lambda,&\lambda'-\mathfrak{L},&\lambda''\\
\beta,&\beta'-\mathfrak{B},&\beta''
\end{array}
\]
zum Grunde legt. Der Erfolg wird sein, dass der nach den daraus folgenden neuen Elementen berechnete Ort von $\lambda'$, $\beta'$ so wenig verschieden ist (bei meinen Erfahrungen nur in Theilen von Secunden), dass es keiner andern Verbesserung bedarf.

\begin{center}9.\end{center}
Das so eben angezeigte Verfahren gilt nur für den Fall, da man die Bahnbestimmung auf \emph{dieselben} Beobachtungen gründen will, die man zur ersten An%
\SourceLeaf{169}%
\expandafter\def\csname HistoricalNoteText12\endcsname{V.\ Theoria motus corporum coelest.\ art.\ 83.}%
näherung angewandt hat. Wenn man aber nachher die Verbesserung der Elemente durch lauter oder zum Theil andere Beobachtungen sucht, so habe ich nach mancherlei andern Proben folgende zwei Methoden am brauchbarsten gefunden.

I. Man berechnet aus den zwei äussern geocentrischen Örtern die heliocentrischen nach 3 Hypothesen, indem man zuerst die genäherten Abstände für diese Beobachtungen voraussetzt, und nachher erst den einen, dann den andern ein wenig ändert. Nach den in allen 3 Hypothesen gefundenen Elementen berechnet man den Ort für die mittlere Beobachtung, den man mit dem beobachteten vergleicht. Durch Interpolation findet man sodann die corrigirten Abstände, und wenn man will, auch die corrigirten Elemente, doch ist es besser die Mühe nicht zu scheuen, diese durch besondere Rechnung aus den neuen Abständen zu berechnen, zumal wenn die Änderungen der Elemente noch sehr stark sind.

II\textsuperscript{a}. Man bedient sich ganz desselben Verfahrens, nur mit dem Unterschiede, dass man statt der genäherten Abstände in den äussern Beobachtungen die genäherte Bestimmung der Neigung und des $\AscendingNode$ gebraucht, und jede von diesen etwas ändert.

II\textsuperscript{b}. Man berechnet theils mit den genäherten theils mit etwas geänderten Bestimmungen der Neigung und des $\AscendingNode$ aus allen drei geocentrischen Örtern die heliocentrischen; aus den zwei äussern heliocentrischen die Elemente und aus diesen den mittlern heliocentrischen, den man mit dem aus dem beobachteten geocentrisch abgeleiteten vergleicht und dann die verbesserte Neigung und $\AscendingNode$ durch Interpolation sucht u.\ s.\ w.

Man könnte auch in II\textsuperscript{b} aus den \emph{drei heliocentrischen} Örtern nach bekannten Formeln die Ellipse bestimmen, ohne die Zeiten mit anzuwenden; aus den Dimensionen der Ellipse die 2 Zwischenzeiten berechnen und mit den wahren vergleichen und dann eben so wie vorher die corrigirte Neigung und $\AscendingNode$ durch Interpolation suchen. Allein dies Verfahren habe ich nach meiner Erfahrung verwerfen müssen. Man würde auf diesem Wege nur nach wiederholten Operationen mit weit mehr Mühe zu einer genauen Darstellung der Beobachtungen gelangen. Die Ursachen davon hier ausführlich zu untersuchen, würde zu weitläufig sein\HistoricalNote{12}. Ich bemerke daher nur, dass man auf diese Art die zweiten Dif%
\SourceLeaf{170}%
ferentiale, von denen man sich gerade durch die im 5.\ und 6.\ Artikel ausgeführten Kunstgriffe losgemacht hat, wieder herbeiführt, und dass diese delicaten zweiten Differentiale durch eine nicht grosse Veränderung der Neigung und des $\AscendingNode$ \emph{ganz enorm} entstellt werden können. zumal wenn die Excentricität nicht gross ist. Es kann hier leicht kommen, dass ein paar Minuten Änderung im $\AscendingNode$ oder der Neigung eine Ellipse hervorbringen kann, die mit der vorhergehenden fast gar keine Ähnlichkeit hat, daher denn begreiflich der Interpolation nicht mehr getraut werden kann. Dies ist nicht der Fall bei unsern Methoden, wo immer nur zwei Beobachtungen zum Grunde gelegt werden. \emph{Sapienti sat.} — Nach meinen wiederholten Erfahrungen finde ich die Methode I am allerzweckmässigsten und allgemeinsten.

Übrigens gelten alle diese Methoden nur, so lange der Bogen noch mässig gross ist. Hat man schon Beobachtungen von 1 oder mehrern Jahren, so werden wieder andere nöthig sein, über die ich mich hier nicht weitläufig ausbreiten kann. In diesem Falle ist es im Allgemeinen nicht anzurathen, die Elemente auf drei vollständige Beobachtungen zu gründen, sondern es ist weit angemessener vier Längen und zwei Breiten zu gebrauchen. — Umfassen die Beobachtungen schon \emph{mehrere} Jahre, und sind die Elemente schon \emph{bis auf Kleinigkeiten} bestimmt, so halte ich den Gebrauch der Differential-Änderungen, wobei man eine beliebige Zahl von Beobachtungen zum Grunde legen kann, für das beste Mittel.
